\begin{table}[t!]
\small
    \centering
  \caption{Experiment Dataset Statistics}
  \vspace{-0.1in}
  \label{tab:stats}
  \begin{tabular}{cccccccccc}
    \toprule
    Dataset & \# Users & \# Items & \# Interactions\\
    \midrule
    MIND \cite{wu2020mind} & 1,000,000 & 160,000 & 15,000,000\\
    KuaiRec \cite{gao2022kuairec} & 1,411 & 3,327 & 4,676,570\\
  \bottomrule
\end{tabular}
% \vspace{-0.3in}
\end{table}

\begin{table*}[h]
\small
  \caption{MIND Dataset Illustration}
  \vspace{-0.1in}
  \label{tab:mind_example}
  \begin{tabular}{ccccc}
    \toprule
    Impression ID & User ID & Time & User Interest History & News with Labels\\
    \midrule 
    91 & U397059 & 11/15/2019 10:22:32 AM & N106403 N71977 N97080 N102132 & N129416-0 N26703-1 N120089-1 N53018-0 N89764-0\\
  \bottomrule
\end{tabular}
\vspace{-0.15in}
\end{table*}

\begin{table}[t!]
\small
  \caption{Parameter Sizes across Different Architectures}
  \vspace{-0.1in}
  \label{tab:param}
  \begin{tabular}{ccccc}
    \toprule
    Algorithm & MIND & KuaiRec\\
    \midrule 
    MLP & $0.76 \times 10^6$ & $1.21 \times 10^7$\\
    Wide \& Deep & $/$ & $2.33 \times 10^7$\\
    NCF & $1.51 \times 10^6$ & $2.31 \times 10^7$\\
    NRMS & $1.49 \times 10^6$ & $/$\\ 
    NAML & $1.49 \times 10^6$ & $/$\\
    DeepFM & $/$ & $2.30 \times 10^7$ \\
    EpiNet + MLP & $0.77 \times 10^6$ & $1.22 \times 10^7$\\
    EpiNet + Wide \& Deep & $/$ & $2.34 \times 10^7$\\
    EpiNet + NCF & $1.52 \times 10^6$ & $2.32 \times 10^7$\\
    EpiNet + NRMS & $1.50 \times 10^6$ & $/$ \\
    EpiNet + NAML & $1.50 \times 10^6$ & $/$\\
    EpiNet + DeepFM & $/$ & $2.31 \times 10^7$ \\
    ENR & $1.52 \times 10^6$ & $2.32 \times 10^7$ \\
  \bottomrule
\end{tabular}
\vspace{-0.1in}
\end{table}

\begin{table}[t!]
\small
  \caption{Average Computation Time for MIND Experiments (Seconds). Both full Neural TS and Neural UCB agents require inverting a covariance matrix of approximate size $10^{6} \times 10^{6}$, which leads to $\infty$ inference time and intractable memory. Even when only using the last layer gradients, Neural TS and Neural UCB still takes about 10 times of ENR agent's inference time.}
  \vspace{-0.1in}
  \label{tab:mind_comp}
  \begin{tabular}{ccccc}
    \toprule
    Algorithm & Inference / Action & Training\\
    \midrule 
    Neural UCB (w/ any arch. below) & $\infty$ & $\infty$\\
    Neural TS (w/ any arch. below) & $\infty$ & $\infty$\\
    \midrule
    $\epsilon$-Greedy + MLP  & $0.0006 \pm 0.0001$ & $0.11 \pm 0.01$\\
    Neural UCB (Last Layer) + MLP & $0.0098 \pm 0.001$ & $0.13 \pm 0.01$\\
    Neural LinUCB + MLP & $0.0045 \pm 0.001$ & $0.13 \pm 0.01$\\
    Neural TS (Last Layer) + MLP & $0.0099 \pm 0.001$ & $0.14 \pm 0.01$\\
    Ensemble + MLP & $0.0006 \pm 0.0001$ & $0.72 \pm 0.03$\\

    \midrule

    $\epsilon$-Greedy + NCF & $0.00079 \pm 0.0001$ & $0.13 \pm 0.01$\\
    Neural UCB (Last Layer) + NCF & $0.014 \pm 0.002$ & $0.16 \pm 0.01$\\
    Neural LinUCB + NCF & $0.0062 \pm 0.001$ & $0.15 \pm 0.01$\\
    Neural TS (Last Layer) + NCF & $0.013 \pm 0.002$ & $0.16 \pm 0.01$\\
    Ensemble + NCF & $0.00081 \pm 0.0001$ & $0.93 \pm 0.03$\\

    \midrule
    
    $\epsilon$-Greedy + NAML & $0.00082 \pm 0.0001$ & $0.13 \pm 0.01$\\
    Neural UCB (Last Layer) + NAML &$0.014 \pm 0.002$ & $0.15 \pm 0.01$\\
    Neural LinUCB + NAML &$0.063 \pm 0.001$ & $0.15 \pm 0.01$\\
    Neural TS (Last Layer) + NAML & $0.014 \pm 0.002$ & $0.15 \pm 0.01$\\
    Ensemble + NAML & $0.00079 \pm 0.0001$ & $0.95 \pm 0.03$\\

    \midrule

  $\epsilon$-Greedy + NRMS & $0.00083 \pm 0.0001$ & $0.14 \pm 0.01$\\
    Neural UCB (Last Layer) + NRMS & $0.013 \pm 0.002$ & $0.16 \pm 0.01$\\
    Neural LinUCB + NRMS & $0.006 \pm 0.001$ & $0.15 \pm 0.01$\\
    Neural TS (Last Layer) + NRMS & $0.014 \pm 0.002$ & $0.15 \pm 0.01$\\
    Ensemble + NRMS & $0.00078 \pm 0.0001$ & $0.91 \pm 0.03$\\

    \midrule\midrule
    EpiNet + MLP & $0.0007 \pm 0.0001$ & $0.12 \pm 0.01$\\
    EpiNet + NCF & $0.0012 \pm 0.0001$ & $0.16 \pm 0.01$\\
    EpiNet + NAML & $0.0013 \pm 0.0001$ & $0.17 \pm 0.01$\\
    EpiNet + NRMS & $0.0012 \pm 0.0001$ & $0.17 \pm 0.01$\\
    \midrule
    ENR & $0.00095 \pm 0.0001$ & $0.18 \pm 0.01$\\
  \bottomrule
\end{tabular}
\vspace{-0.2in}
\end{table}

\begin{table}[t!]
\small
  \caption{Average Computation Time for KuaiRec Experiments (Seconds). Both full Neural TS and Neural UCB agents require inverting a covariance matrix of approximate size $10^{7} \times 10^{7}$, which leads to $\infty$ inference time and intractable memory. Note that NeuralUCB and NeuralTS agents that only compute the last layer of gradients still takes 150 times of ENR's inference time due to available action size of 3327 for every time step, hence intractable in this experiment.}
  \vspace{-0.1in}
  \label{tab:kuai_comp}
  \begin{tabular}{ccccc}
    \toprule
    Algorithm & Inference / Action & Training\\
    \midrule
    Neural UCB (w/ any arch. below) & $\infty$ & $\infty$\\
    Neural TS (w/ any arch. below) & $\infty$ & $\infty$\\
    \midrule
    $\epsilon$-Greedy + MLP & $0.012 \pm 0.001$ & $0.24 \pm 0.01$\\
    Neural UCB (Last Layer) + MLP & $1.60 \pm 0.13$ & $0.25 \pm 0.01$\\
    Neural LinUCB + MLP & $0.036 \pm 0.002$ & $0.25 \pm 0.01$\\
    Neural TS (Last Layer) + MLP & $1.59 \pm 0.13$ & $0.25 \pm 0.01$\\
    Ensemble + MLP & $0.018 \pm 0.001$ & $1.2 \pm 0.03$\\
    \midrule
    $\epsilon$-Greedy + NCF & $0.013 \pm 0.001$ & $0.31 \pm 0.02$\\
    Neural UCB (Last Layer) + NCF & $1.71 \pm 0.13$ & $0.35 \pm 0.02$\\
    Neural LinUCB + NCF & $0.042 \pm 0.002$ & $0.36 \pm 0.02$\\
    Neural TS (Last Layer) + NCF & $1.72 \pm 0.13$ & $0.36 \pm 0.02$\\
    Ensemble + NCF & $0.020 \pm 0.001$ & $1.5 \pm 0.03$\\
    \midrule
    $\epsilon$-Greedy + Wide \& Deep & $0.013 \pm 0.001$ & $0.32 \pm 0.02$\\
    Neural UCB (Last Layer) + Wide \& Deep & $1.68 \pm 0.13$ & $0.34 \pm 0.02$\\
    Neural LinUCB + Wide \& Deep & $0.039 \pm 0.002$ & $0.35 \pm 0.02$\\
    Neural TS (Last Layer) + Wide \& Deep & $1.69 \pm 0.13$ & $0.34 \pm 0.02$\\
    Ensemble + Wide \& Deep & $0.020 \pm 0.001$ & $1.6 \pm 0.03$\\
    
    \midrule
    $\epsilon$-Greedy + DeepFM & $0.014 \pm 0.001$ & $0.32 \pm 0.02$\\
    Neural UCB (Last Layer) + DeepFM & $1.72 \pm 0.13$ & $0.36 \pm 0.02$\\
    Neural LinUCB + DeepFM & $0.041 \pm 0.002$ & $0.36 \pm 0.02$\\
    Neural TS (Last Layer) + DeepFM & $1.72 \pm 0.13$ & $0.36 \pm 0.02$\\
    Ensemble + DeepFM & $0.021 \pm 0.001$ & $1.6 \pm 0.03$\\
    \midrule\midrule
    EpiNet + MLP & $0.013 \pm 0.001$ & $0.28 \pm 0.01$\\
    EpiNet + NCF & $0.015 \pm 0.001$ & $0.37 \pm 0.01$\\
    EpiNet + Wide \& Deep & $0.015 \pm 0.001$ & $0.38 \pm 0.01$\\
    EpiNet + DeepFM & $0.014 \pm 0.001$ & $0.38 \pm 0.01$\\
    \midrule
    ENR & $0.012 \pm 0.001$ & $0.38 \pm 0.01$\\
  \bottomrule
\end{tabular}
\vspace{-0.1in}
\end{table}

In this section, we will go through a set of experiments ranging from a toy experiment to large-scale real-world data experiments. The dataset statistics are shown in Table \ref{tab:stats}. Each of the dataset has millions of interactions. We will compare ENR and EpiNet \cite{osband2021epistemic} (first time adopted for neural contextual bandits) against $\epsilon$-Greedy, Ensemble Sampling with Prior Networks (Ensemble) \cite{lu2017ensemble}, Neural Thompson Sampling (Neural TS) \cite{zhang2021neural}, Neural UCB \cite{zhou2020neural}, and Neural LinUCB \cite{xu2021neural}, each with a similar sized neural network. In addition, to assess advantages attributable to the architecture of ENR, we will also evaluate it against a few well recognized RS neural network architecture including, Neural Collaborative Filtering (NCF) \cite{he2017neural}, Neural Recommendation with Attentive Multi-View Learning (NAML) \cite{wu2019naml}, Neural Recommendation with Multi-Head Self-Attention (NRMS) \cite{wu2019neural}, Wide and Deep \cite{cheng2016wide}, and DeepFM \cite{guo2017deepfm}, combined with exploration strategies above. To ensure fairness of evaluation, we ensure that all neural network architectures other than vanilla MLP has a similar parameter size. See Table \ref{tab:param}. The algorithms all adopt ADAM \cite{kingma2015adam} for model optimization. Different from EpiNet, ENR empirically requires additional L2 regularization in addition to prior network regularization. We choose L2 decay coefficient to be 0.0001 across the experiments. 

Note that for both Neural UCB and Neural Thompson Sampling, due to their scalability, we are only able to experiment them with their last layer versions, where we only compute gradients of the last layer of the neural networks. The full Neural UCB and Neural Thompson Sampling agents require inverting matrices of sizes $10^6 \times 10^6$ (7.45 TB of memory and $O(10^{18})$ complexity) and $10^7 \times 10^7$ (745 TB of memory and $O(10^{21})$ complexity) for MIND and KuaiRec experiments respectively, which are intractable in most commercial machines. See more details of computation costs in Table \ref{tab:mind_comp} and \ref{tab:kuai_comp} respectively for two experiments and all experiments are performed on an AWS p4d24xlarge machine with 8 A100 GPUs, each GPU with 40 GB of GPU memory. NeuralLinUCB requires 6x and 3.5x inference time compared to our method in two real-world experiments with a similar training time.  Simplified versions of NeuralTS and NeuralUCB using only gradients from the last layer of NNs require 10x and 100x inference time compared to our method in these experiments. Lastly, Ensemble requires a similar inference time, but requires 5x training time compared to our method with parallel computing optimization. NeuralLinUCB, NeuralUCB, NeuralTS all require orders of magnitude higher inference cost and Ensemble requires orders of magnitude higher training cost.

We will use average reward as performance metric in the following experiments (average click-through rate and average user rating in MIND and KuaiRec respectively). We choose to not use classic metrics like NDCG or precision because they are supervised-learning metrics. In online bandits, the agent’s optimal strategy could try to explore and retrieve information instead of optimizing these metrics.  

\begin{table*}[t!]
\small
  \caption{Average Regret for Toy Experiments}
  \vspace{-0.1in}
  \label{tab:toy_regret}
  \begin{tabular}{ccccccccccc}
    \toprule
    Algorithm & Exploit & $\epsilon$-Greedy & Neural UCB & Neural LinUCB & Neural TS & Ensemble & EpiNet & ENR\\
    \midrule 
    Avg Regret & $0.10 \pm 0.006$ & $0.11 \pm 0.01$ & $0.11 \pm 0.007$ & $0.07 \pm 0.009$ & $0.075 \pm 0.008$ & $0.08 \pm 0.004$ & $0.064 \pm 0.003$ & $\mathbf{0.053 \pm 0.0006}$\\
  \bottomrule
\end{tabular}
\vspace{-0.1in}
\end{table*}

\begin{table*}[t!]
\small
  \caption{User Interactions (Thousands) Needed for Best Algorithm Candidates (Selected among Different Neural Network Architectures) to Reach Standard Performance and Savings Compared to Non-EpiNet Baselines}
  \vspace{-0.1in}
  \label{tab:sample_complexity}
  \begin{tabular}{c|cccccc|cc|cc}
    \toprule
    Algorithm & Exploit & $\epsilon$-Greedy & Neural UCB & Neural LinUCB & Neural TS & Ensemble & EpiNet & ENR & Savings\\
    \midrule 
    MIND ($0.2$ Click-Through Rate) & $\infty$ & $\infty$ & $277$ & $268$ & $184$ & $142$ & $145$ & \textbf{62} & $56\%$\\
    KuaiRec ($2.3$ User Rating) & N/A & $4.5$ & N/A & $4$ & N/A & $2.4$ & $2.1$ & \textbf{1.7} & $29\%$\\
  \bottomrule
\end{tabular}
\vspace{-0.1in}
\end{table*}

 \begin{figure}
     \centering
     \includegraphics[width=0.3\textwidth]{img/toy_1.png}
     \caption{Toy Experiment Results}
     \label{fig:toy_1}
     \vspace{-0.1in}
 \end{figure}

\begin{table}[t!]
\small
  \caption{Average Click-Through Rate for MIND Experiments}
  \vspace{-0.1in}
  \label{tab:mind_reward}
  \begin{tabular}{ccccc}
    \toprule
    Algorithm & MIND Train & MIND Eval\\
    \midrule 
    Exploit + MLP & $0.170 \pm 0.0007$ & $0.168 \pm 0.0011$ \\
    $\epsilon$-Greedy + MLP  & $0.171 \pm 0.0007$ & $0.167 \pm 0.0011$\\
    Neural UCB (Last Layer) + MLP & $0.188 \pm 0.0016$ & $0.192 \pm 0.0010$\\
    Neural LinUCB + MLP & $0.189 \pm 0.0032$ & $0.192 \pm 0.00086$\\
    Neural TS (Last Layer) + MLP & $0.191 \pm 0.0023$ & $0.194 \pm 0.0015$\\
    Ensemble + MLP & $0.193 \pm 0.0022$ & $0.191 \pm 0.0009$\\

    \midrule
    Exploit + NCF & $0.181 \pm 0.0022$ & $0.179 \pm 0.0024$\\
    $\epsilon$-Greedy + NCF & $0.182 \pm 0.0023$ & $0.178 \pm 0.0025$\\
    Neural UCB (Last Layer) + NCF& $0.203 \pm 0.0029$ & $0.203 \pm 0.0023$\\
    Neural LinUCB + NCF & $0.203 \pm 0.0029$ & $0.203 \pm 0.0030$\\
    Neural TS (Last Layer) + NCF & $0.207 \pm 0.0031$ & $0.206 \pm 0.0028$\\
    Ensemble + NCF & $0.209 \pm 0.0040$ & $0.206 \pm 0.0031$\\
    
    \midrule
    Exploit + NAML & $0.180 \pm 0.0021$ & $0.178 \pm 0.0024$\\
    $\epsilon$-Greedy + NAML & $0.182 \pm 0.0022$ & $0.178 \pm 0.0025$\\
    Neural UCB (Last Layer) + NAML & $0.202 \pm 0.0028$ & $0.203 \pm 0.0024$\\
    Neural LinUCB + NAML & $0.203 \pm 0.0028$ & $0.203 \pm 0.0030$ \\
    Neural TS (Last Layer) + NAML & $0.206 \pm 0.003$ & $0.207 \pm 0.0028$\\
    Ensemble + NAML & $0.208 \pm 0.0039$ & $0.206 \pm 0.0031$\\

    \midrule
    Exploit + NRMS & $0.181 \pm 0.0022$ & $0.178 \pm 0.0024$\\
   $\epsilon$-Greedy + NRMS & $0.182 \pm 0.0023$ & $0.178 \pm 0.0025$\\
    Neural UCB (Last Layer) + NRMS & $0.203 \pm 0.0029$ & $0.203 \pm 0.0024$\\
    Neural LinUCB + NRMS & $0.204 \pm 0.0029$ & $0.203 \pm 0.0030$\\
    Neural TS (Last Layer) + NRMS & $0.206 \pm 0.0030$ & $0.207 \pm 0.0027$\\
    Ensemble + NRMS & $0.208 \pm 0.0039$ & $0.206 \pm 0.0031$\\

    \midrule\midrule
    EpiNet + MLP & $0.204 \pm 0.003$ & $0.199 \pm 0.0012$\\
    EpiNet + NCF & $0.219 \pm 0.004$ & $0.214 \pm 0.0028$\\
    EpiNet + NAML & $0.217 \pm 0.004$ & $0.214 \pm 0.0028$\\
    EpiNet + NRMS & $0.218 \pm 0.004$ & $0.214 \pm 0.0027$\\
    \midrule
    ENR & $\mathbf{0.228 \pm 0.004}$ & $\mathbf{0.220 \pm 0.0026}$\\
  \bottomrule
\end{tabular}
\vspace{-0.15in}
\end{table}

\begin{table}[t!]
\small
  \caption{Average Reward for KuaiRec Experiments}
  \vspace{-0.1in}
  \label{tab:kuai_reward}
  \begin{tabular}{ccccc}
    \toprule
    Algorithm & KuaiRec Train & KuaiRec Eval\\
    \midrule 
    $\epsilon$-Greedy + MLP & $1.41 \pm 0.087$ & $1.57 \pm 0.11$\\
    Neural LinUCB + MLP & $1.68 \pm 0.15$ & $1.59 \pm 0.13$\\
    Ensemble + MLP & $2.28 \pm 0.061$ & $2.03 \pm 0.09$\\
    \midrule
    $\epsilon$-Greedy + NCF & $2.42 \pm 0.0075$ & $2.31 \pm 0.019$\\
    Neural LinUCB + NCF & $2.47 \pm 0.022$ & $2.38 \pm 0.020$\\
    Ensemble + NCF & $2.55 \pm 0.003$ & $2.46 \pm 0.005$\\
    \midrule
    $\epsilon$-Greedy + Wide \& Deep & $2.38 \pm 0.007$ & $2.27 \pm 0.018$\\
    Neural LinUCB + Wide \& Deep & $2.42 \pm 0.021$ & $2.35 \pm 0.017$\\
    Ensemble + Wide \& Deep & $2.52 \pm 0.003$ & $2.43 \pm 0.005$\\
    \midrule
    $\epsilon$-Greedy + DeepFM & $2.38 \pm 0.0075$ & $2.32 \pm 0.022$\\
    Neural LinUCB + DeepFM & $2.45 \pm 0.022$ & $2.36 \pm 0.019$\\
    Ensemble + DeepFM & $2.53 \pm 0.003$ & $2.47 \pm 0.005$\\
    \midrule\midrule
    EpiNet + MLP & $2.43 \pm 0.01$ & $2.17 \pm 0.087$\\
    EpiNet + NCF & $2.58 \pm 0.0047$ & $2.43 \pm 0.016$\\
    EpiNet + Wide \& Deep & $2.56 \pm 0.005$ & $2.41 \pm 0.018$\\
    EpiNet + DeepFM & $2.57 \pm 0.0046$ & $2.46 \pm 0.021$\\
    \midrule
    ENR & $\mathbf{2.62 \pm 0.0005}$ & $\mathbf{2.63 \pm 0.066}$\\
  \bottomrule
\end{tabular}
\vspace{-0.2in}
\end{table}

% \begin{table*}[t!]
% \small
%   \caption{Average Reward for Real-World Experiments}
%   \vspace{-0.1in}
%   \label{tab:exp_reward}
%   \begin{tabular}{ccccc}
%     \toprule
%     Algorithm & MIND Dataset Train & MIND Dataset Eval & KuaiRec Dataset Train & KuaiRec Dataset Eval\\
%     \midrule 
%     $\epsilon$-Greedy + MLP  & $0.227 \pm 0.009$ & $1.80 \pm 0.48$\\
%     Neural UCB (Last Layer) + MLP & $0.223 \pm 0.022$ & $9.86 \pm 1.11$ & $118.93 \pm 1.65$ \\
%     Neural LinUCB + MLP & $0.215 \pm 0.043$ & & \\
%     Neural TS (Last Layer) + MLP & $0.222 \pm 0.011$ & $6.55 \pm 1.14$ & $114.97 \pm 1.59$ \\
%     Ensemble + MLP & $0.227 \pm 0.034$ & $2.51 \pm 0.02$\\
%     $\epsilon$-Greedy + NCF & $0.249 \pm 0.039$ & $2.46 \pm 0.04$\\
%     % Neural UCB & $0.0 \pm 0.0$ & $9.86 \pm 1.11$ & $118.93 \pm 1.65$ \\
%     % Neural TS & $0.0 \pm 0.0$ & $6.55 \pm 1.14$ & $114.97 \pm 1.59$ \\
%     % Ensemble + NCF & $0.26 \pm 0.01$ & $2.56 \pm 0.03$\\
%     \midrule
%     EpiNet + MLP & $0.269 \pm 0.016$ & $2.45 \pm 0.05$\\
%     % EpiNet + NCF & $0.32 \pm 0.01$ & $2.60 \pm 0.01$\\
%     \midrule
%     ENR & $\mathbf{0.285 \pm 0.017}$ & $\mathbf{2.64 \pm 0.00}$\\
%   \bottomrule
% \end{tabular}
% \end{table*}

% \begin{table*}[ht!]
% \small
%     \centering
%   \caption{KuaiRec Experiment Algorithm Prediction \& Training Computation Time with MLP Architecture (Seconds)}
%   \vspace{-0.1in}
%   \label{tab:comp_time}
%   \begin{tabular}{cccccccccc}
%     \toprule
%     Task & $\epsilon$-Greedy & Ensemble & EpiNet & \\
%     \midrule
%     Pred & 0.063 & 0.065 & 0.07 & $\infty$ & $\infty$ & 0.24 & 0.21 \\
%     Train & 0.56 & 1.34 & 0.82 & 0.53 & 0.64 & 0.53 & 0.64\\
%   \bottomrule
% \end{tabular}
% \end{table*}


\subsection{Toy Experiment}
For this experiment, we define a synthetic environment to evaluate the performance of ENR compared to various baselines mentioned above. According to Section \ref{sec:prob}, the synthetic environment is characterized as: Both action and context vectors are of dimension 100. $R_{t+1} \sim \text{Bernoulli}\left(\sigma(\theta^T\text{concat}(\psi_{S_t}, \phi_{A_t}))\right)$. The size of $\mathcal{A}$ is 100. $\sigma$ is the Sigmoid function and both $\theta$ is unknown to all the agents as the learning target. All of $\Theta, \phi_A, \psi_S$ are sampled from standard Gaussian. Since the toy environment is synthesized, we could compute the regret of the algorithms as
\begin{equation}
    \text{Regret}(T) = \frac{1}{T}\sum_{t=1}^T (R_{t, A^*_t} - R_{t, A_t}),
\end{equation}
where $R_{t, A^*_t}$ indicates the reward of the optimal action given the context $S_t$. In this experiment, we use a hidden layer of (500, 100, 50) for MLP and the main neural network architecture for MLP and ENR. The action and context summarization layer for ENR is set without hidden layers and with activation of ReLU with Layer Normalization on top of that. The uncertainty estimation architecture of ENR is set with hidden layer of (100, 50). The EpiNet architecture is set to a single hidden layer of 50. All prior network scaling is set to 0.3. Dimension of epistemic indices is set to 5 for EpiNet and ENR. The results are averaged over 50 independent experiments. The experiment is run for 5000 steps.

See Figure \ref{fig:toy_1} and Table \ref{tab:toy_regret} for more details of the experiment results. We see that ENR marginally improves on top of EpiNet with a similarly sized base neural network of fully connected layers. Given that the setups of both toy experiments are relatively simple, it is expected that sophisticated neural representations would not make material difference.  

\begin{figure*}[ht!]
     \begin{subfigure}[t]{0.24\textwidth}
         \centering
         \includegraphics[width=\textwidth]{img/mind_mlp_train.png}
         \caption{MIND Experiment MLP Training Results}
         \label{fig:mind_mlp}
         \vspace{-0.15in}
     \end{subfigure}
     \hfill
     \begin{subfigure}[t]{0.24\textwidth}
         \centering
         \includegraphics[width=\textwidth]{img/mind_ncf_train.png}
         \caption{MIND Experiment NCF Training Results}
         \label{fig:mind_ncf}
         \vspace{-0.15in}
     \end{subfigure}
     \hfill
     \begin{subfigure}[t]{0.24\textwidth}
         \centering
         \includegraphics[width=\textwidth]{img/mind_naml_train.png}
         \caption{MIND Experiment NAML Training Results}
         \label{fig:mind_naml}
         \vspace{-0.15in}
     \end{subfigure}
     \hfill
     \begin{subfigure}[t]{0.24\textwidth}
         \centering
         \includegraphics[width=\textwidth]{img/mind_nrms_train.png}
         \caption{MIND Experiment NRMS Training Results}
         \label{fig:mind_nrms}
         \vspace{-0.15in}
     \end{subfigure}
     \caption{MIND Experiments Contextual Bandit Training Outcome}
     \label{fig:mind}
     \vspace{-0.1in}
\end{figure*}

\begin{figure*}[t!]
     \begin{subfigure}[t]{0.24\textwidth}
         \centering
         \includegraphics[width=\textwidth]{img/kuai_mlp_train.png}
         \caption{KuaiRec Experiment MLP Training Results}
         \label{fig:mind_mlp}
         \vspace{-0.15in}
     \end{subfigure}
     \hfill
     \begin{subfigure}[t]{0.24\textwidth}
         \centering
         \includegraphics[width=\textwidth]{img/kuai_ncf_train.png}
         \caption{KuaiRec Experiment NCF Training Results}
         \label{fig:mind_ncf}
         \vspace{-0.15in}
     \end{subfigure}
     \hfill
     \begin{subfigure}[t]{0.24\textwidth}
         \centering
         \includegraphics[width=\textwidth]{img/kuai_wd_train.png}
         \caption{KuaiRec Experiment Wide \& Deep Training Results}
         \label{fig:mind_wd}
         \vspace{-0.15in}
     \end{subfigure}
     \hfill
     \begin{subfigure}[t]{0.24\textwidth}
         \centering
         \includegraphics[width=\textwidth]{img/kuai_deepfm_train.png}
         \caption{KuaiRec Experiment DeepFM Training Results}
         \label{fig:mind_deepfm}
         \vspace{-0.15in}
     \end{subfigure}
     \caption{KuaiRec Experiments Contextual Bandit Training Outcome}
     \label{fig:kuai}
     \vspace{-0.1in}
\end{figure*}

\subsection{MIND Dataset Experiment \cite{wu2020mind}}
MIND is a large-scale dataset generated by Microsoft news recommender system. It provides anonymized data logs that record each user's historical interesting articles and their impression as well as click logs with new recommendations. The dataset contains over millions of anonymized interaction logs. For statistics of the dataset, please refer to Table \ref{tab:stats}. One unique advantage of the MIND dataset for contextual bandit setup is its complete ground truth. At time step $t$, the MIND dataset logs $K$ distinct news recommendations sent to the user and records the news articles the user gave good feedback on, all of which happen at a single timestamp. Please see Table \ref{tab:mind_example} for an illustration of the data we use.  Hence when evaluating the outcome of any action chosen by an agent among the sent recommendations, we can directly use the groundtruth from the dataset and avoid counterfactual evaluation which leads to high variance. 

 In this experiment, we use all exploration baseline strategies as well as NCF \cite{he2017neural}, NAML \cite{wu2019naml} and NRMS \cite{wu2019neural} for neural network architecture baselines. It is worth noting that in this experiment, we adopt the Multi-Head Self-Attention architecture mentioned in NRMS \cite{wu2019neural} to extract representation for context (user) before feeding into the neural networks (except NAML where it has its own summarization module). We did not select Wide \& Deep and DeepFM because both neural architectures are heavily built on sparse feature optimization, which is not present in this experiment.

In the following, we use a hidden layer of (1000, 500, 100, 50) for MLP and the final neural network architecture for NCF and ENR. NCF's first MLP uses hidden layer of (1000, 500). For NAML and NRMS, EpiNet is added to the last layer replacing dot product with item-wise multiplication. The action and context summarization layer for ENR is set without hidden layers and with activation of ReLU and Layer Normalization on top of that. The uncertainty estimation architecture of ENR is set with hidden layer of (1000, 500). The EpiNet architecture for other neural networks are set with a single hidden layer of 50. All prior network scaling is set to 0.3. Dimension of epistemic indices is set to 5 for all EpiNets and ENR. The results are averaged over 10 independent experiments. Each training experiment is run for 375,000 time steps by sampling 375,000 users randomly from the dataset. In addition, we also sample 2,000 users that are never seen by the agent as evaluation set for out of sample evaluation.

See Figure \ref{fig:mind} as well as Table \ref{tab:mind_reward} for the experiment results. ENR is able to outperform all other baselines with a big margin in both online interaction and offline out-of-sample evaluation, outperforming $9\%$ and $6\%$ respectively in Click-Through Rate compared to the best baseline candidate (excluding EpiNet candidates, as that is considered part of the contribution of this work) respectively. Moreover, ENR is able to achieve a standard performance of $0.2$ click-through rate with $56\%$ less user interactions compared to the best non-EpiNet baseline. See Table \ref{tab:sample_complexity} for more details.

\begin{figure*}[t!]
     \begin{subfigure}[t]{0.3\textwidth}
         \centering
         \includegraphics[width=\textwidth]{img/enr_layernorm.png}
         \caption{Comparison between with and without Layer Normalization Based on KuaiRec}
         \label{fig:layernorm}
         \vspace{-0.15in}
     \end{subfigure}
     \hfill
     \begin{subfigure}[t]{0.3\textwidth}
         \centering
         \includegraphics[width=\textwidth]{img/enr_index_width.png}
         \caption{Selection of Index Width Based on KuaiRec}
         \label{fig:width}
         \vspace{-0.15in}
     \end{subfigure}
     \hfill
     \begin{subfigure}[t]{0.3\textwidth}
         \centering
         \includegraphics[width=\textwidth]{img/enr_prior.png}
         \caption{Selection of Prior Scale Based on KuaiRec}
         \label{fig:scale}
         \vspace{-0.15in}
     \end{subfigure}
     \caption{Ablation Studies}
     \label{fig:ablation}
     \vspace{-0.2in}
\end{figure*}

\vspace{-10pt}
\subsection{KuaiRec Dataset Experiment \cite{gao2022kuairec}}
The last experiment we run is through the KuaiRec dataset. KuaiRec offers almost a full dense matrix in terms of the interactions between users and actions. The density is at 99.6\%. For more details of the dataset statistics, please refer to Table \ref{tab:stats}. With these characteristics of the dataset, we can evaluate all of the algorithms in real-world contextual bandit environment with huge available set of recommendations at every step without resorting to counterfactual evaluation. It is worth noting that, different from MIND, KuaiRec's user ratings on each recommendation are collected at slightly different times, and hence present a less rigorous grounding for avoiding counterfactual evaluation. Nevertheless, we believe it is valuable to still study the performance of ENR and EpiNet algorithms when facing a huge available set of actions at every timestep. We also note that such datasets, with feedback collected at different times, are commonly used in bandit learning literature to evaluate algorithms' performance, e.g. \cite{chris2018learning, hong2020latent, zhang2020conversational}.

In this experiment, we adopt all exploration strategies except Neural Thompson Sampling and Neural UCB due to their computation complexity. Given that both strategies requires a gradient pass on every action from the candidate list, these algorithms do not scale to a problem with 3,327 actions at every time step. We also choose NCF, Wide \& Deep and DeepFM as neural network architectures given that the features provided in the dataset are raw features with many categorical features. The dataset does not provide user historical articles and hence NAML and NRMS are not good fits.

In the following, we use a hidden layer of (4096, 1024, 512, 128) for MLP and the final neural network architecture for NCF and ENR. NCF's first MLP uses hidden layer of (1024, 512). Wide \& Deep and DeepFM adopts the same hidden layer structure for dense embeddings. The action and context summarization layer for ENR is set without hidden layers and with activation of ReLU with Layer Normalization on top of that. The uncertainty estimation architecture of ENR is set with hidden layer of (1024, 128). The EpiNet architecture for other neural networks are set with a single hidden layer of 128. All prior network scaling is set to 0.3. Dimension of epistemic indices is set to 5 for all EpiNets and ENR. The results are averaged over 10 random seeds. Each experiment is run for 6,000 time steps where each context is randomly sampled from the small matrix excluding an evaluation set of 100 contexts. One unique thing to note in this dataset is that user feedback for items in this dataset are rating scalars. Therefore, we use regression and mean squared error loss instead of cross entropy loss.

See Figure \ref{fig:kuai} as well as Table \ref{tab:kuai_reward} to see the detailed results of the experiment. From the results, we can see that ENR outperforms all other candidates in both in-sample online interactions as well as offline out-of-sample evaluation set by $4\%$ and $6\%$ respectively. Another observation we draw from the result is that with a complete action space provided to the agent at every timestep, all algorithms significantly improve in generalization capabilities from in-sample predictions to out-of-sample predictions. In Table \ref{tab:sample_complexity}, we also show that ENR requires $29\%$ less user interactions compared to the best non-EpiNet baseline to achieve a standard performance of $2.3$ average user rating. 

\vspace{-0.1in}
\subsection{Ablation Studies}
In this section, we cover some ablation studies regarding some of our hyperparameter choices as well as some design choices.
\vspace{-0.05in}
\subsubsection{The Importance of Layer Normalization}
In ENR, one of a key components for action and context summarization is Layer Normalization. As ENR is designed to take in any context and action inputs, it is common to observe a few features that are outsized. To ensure that context action interaction is not dominated by a few particular features, before the item-wise multiplication, we add a Layer Normalization operator to ensure more stability in the neural network as well as getting less impacted by outliers during training. See Figure \ref{fig:layernorm} for a performance comparison between ENR with and without Layer Normalization.
\vspace{-0.05in}
\subsubsection{Selection of Epistemic Index Width}
Epistemic index $z$ is the key to identify a posterior sample from an ENR, and we study the impact of index dimension in this study. See Figure \ref{fig:width}. As the width of the index grows, the performance of the algorithm relatively decays. From the figure, we can also see that index width of 5 outperforms others. Hence we choose 5 for all of our experiments.
\vspace{-0.05in}
\subsubsection{Selection of ENR Uncertainty Prior Network Scale} 
The prior network is a key means to regularize Epistemic Neural Networks by pulling the posterior distribution back from diverging too far away from the prior distribution to avoid distribution collapse. We study the prior network scale for ENR through a sweep. Note that prior scale is a scalar ranging $(0, 1)$. See Figure \ref{fig:scale}. We see that setting prior scale to 0.3 offers the best performance.
